\documentclass{ximera}
\title{Runtime Test 13: Sage String Answers}
\begin{document}
\begin{abstract}
Regression coverage for Sage-generated string values used as canonical answers,
including both ordinary expression answers and explicit string-format answers.
\end{abstract}
\maketitle

\begin{sagesilent}
sage_text_answer = "DNE"
sage_second_text_answer = "NONE"
sage_numeric_answer = 3/4
\end{sagesilent}

\section{Sage string answer normalization}

The following two boxes use the same Sage string.  In both cases the student
should be able to enter \texttt{DNE}; the generated display LaTeX
\verb|\text{\texttt{DNE}}| must not become the canonical answer value.

\begin{problem}
Ordinary answer format:
\[
\answer{\sage{sage_text_answer}}
\]
Typing \texttt{DNE} should be accepted.
\end{problem}

\begin{problem}
Explicit string answer format:
\[
\answer[format=string]{\sage{sage_text_answer}}
\]
Typing \texttt{DNE} should be accepted.  The student should not need to type
any TeX wrapper.
\end{problem}

\begin{problem}
Second arbitrary string value:
\[
\answer[format=string]{\sage{sage_second_text_answer}}
\]
Typing \texttt{NONE} should be accepted.  This guards against a fix that is
special-cased to the literal value \texttt{DNE}.
\end{problem}

\begin{problem}
Numeric Sage control:
\[
\answer{\sage{sage_numeric_answer}}
\]
Typing \texttt{3/4} should still be accepted.
\end{problem}

\end{document}
